\par\bigskip\noindent\textbf{K101 supporting proof: The four strict-source surfaces and their normal bundles}\par\nobreak
This calculation uses the actual source and normal maps of EI, IA,
FO and FL.8. It computes the four strict-source surface terms and
their prescribed finite-quotient corrections. The exceptional and
refinement surfaces and the moving selected reference cap remain
in the full ledger. No total source return is assigned here.

\par

\medskip\noindent\textbf{SS.1. The actual orbifold source and its section lines}

\par\nobreak

On the complementary central source the equation is
\[
 \bar Y^2=\bar X^3-2\chi^2\bar X^2+\bar X,
 \qquad
 \gamma(\chi,\bar X,\bar Y)=(i\chi,-\bar X,-i\bar Y).
 \tag{SS.1}
\]
Its overlap with the native finite source is the literal map
\[
 t_e=\chi^{-4},\qquad x_e=\chi^{-2}\bar X,
 \qquad y_e=\chi^{-3}\bar Y.
 \tag{SS.2}
\]
The native source is \(y_e^2=x_e^3-2x_e^2+t_ex_e\),
with its actual I2 resolution at zero. At O use
\(\eta_e=-x_e/y_e\), \(\eta=-\bar X/\bar Y\).
At T use the resolved coordinate \(b_0=y_e/t_e\)
and \(\bar Y\). The two complete source-normal arrows are
\[
 \eta_e=\chi\eta,
 \qquad b_0=\chi\bar Y.
 \tag{SS.3}
\]
They extend through the native I2 chart by IA.5-\/-10.
In both cases the source coordinate is a regular local
coordinate along the whole section. The two sections are
disjoint: they are O and the distinct two-torsion point T
on each smooth fibre, distinct smooth points on the native
I2 resolution, and distinct points on the good fibre at
\(\chi=0\). These same charts give their actual source
normal bundles.

Let \(\bar S\) denote this proper source with the actual
fourth-root finite charts at infinity. Let L be the pullback
of the base line with frame transition
\[
 e_\infty=\chi^{-1}e_{\rm fin}.
 \tag{SS.4}
\]
This is the genuine fourth-root line: its fourth power has
transition \(q_\infty^{-1}\), where
\(q_\infty=\chi^4\). Thus \(L^4=\mathcal O(F)\)
for a full source fibre F, by this explicit frame map.
On a section its degree is one quarter, including the
order-four finite chart factor. Two different fibres are
disjoint, so \(L^2=0\).

SS.3 makes the coefficient of the source normal line
transform by \(\chi^{-1}\) from the finite to the root
chart. Its frame therefore transforms by \(\chi\).
The normal degree of either section is minus one quarter.
The specified source and base maps consequently give
\[
 O^2=T^2=-\frac14,\qquad O\cdot T=0,
 \qquad L\cdot O=L\cdot T=\frac14,\qquad L^2=0.
 \tag{SS.5}
\]
The point-line frame used throughout is the actual
\(\sigma_a/r_a\). At \(\sigma=0\), both admitted even
a satisfy \(aG=O\), and the point-line trivialization is
literal: in its O chart \(r_a=(\tau-\tau(aG))/\tau\),
and away from those two sections it is the canonical point
section. On the selected Tate disc the chosen representative
has \(P_G^2=Q^h\); its unit valuation forces h zero and
its prescribed limit gives \(P_G=-1\). Thus the full
point/theta comparison restricts to one, without a source
period reindex. These are the precise line maps on the
selected and finite-native overlaps of SS.2.

\par

\medskip\noindent\textbf{SS.2. The common-refinement normal bundle}

\par\nobreak

The physical weights are positive on the line-chain X
coordinate and negative on Y. In an ordinary chain
\(X_jY_j=\sigma\), the strict source graph has
\[
 X_j=\eta s,\quad Y_j=t,\quad v=st,
 \quad\sigma=\eta st.
\]
The native chain has \(A_jB_j=\epsilon\),
\(A_j=\eta_e s_e\), \(B_j=t_e\). Here the symbols
t and \(t_e\) denote normal coefficients in these
two displays. The source-base coordinate is unchanged.
The full chain arrow is
\[
 X_j=\chi^jA_j,\quad Y_j=\chi^{-j-1}B_j,
 \qquad s=\chi^{j+1}s_e,\quad t=\chi^{-j-1}t_e.
 \tag{SS.6}
\]

The actual common38 fan further refines this raw parent.
For the strict face \(C_j,C_{j+1}\), its third rays are
\((1,3,4),(1,2,6),(1,1,6),(1,0,4)\) for j zero
through three. If this ray is \((1,k,l)\), the
complete coordinate map is
\[
 \eta=a,\quad v=a^k n_-n_+,
 \quad V=a^l n_-^j n_+^{j+1}.
\]
Substituting into the actual alternating parent gives
\[
 s_{\rm old}=\frac{V}{\eta^{j+1}v^j}
             =a^{3-j}n_+,
 \qquad t_{\rm old}=n_-.
 \tag{SS.7}
\]
This proves the additional section-ideal power at both O
and T. Away from those sections the factor is invertible
and is the same full parent map. Combining SS.3, SS.6
and SS.7 gives the actual orbifold normal extensions
\[
 N_{+,j}=\mathcal I_{O+T}^{\,4-j}\otimes L^j,
 \qquad N_{-,j}=L^{-j-1}.
 \tag{SS.8}
\]
Their coefficients in the refined O/T root chart have
characters \(\chi^4\) and \(\chi^{-j-1}\).
The section ideal itself supplies \(4-j\) of the
positive powers, and L supplies the remaining j.
This specifies both the line morphism and its finite
character. The full point-line factors have the same
central trivialization described after SS.5; their
off-component units remain in the normal chart maps.

On the selected source chart \(xy=q\), the j-th
sigma-factor chart has coordinates
\[
 \sigma=st,\qquad U=\sigma^j s,
 \qquad V=t\sigma^{3-j}xy,
 \quad0\le j\le3.
 \tag{SS.9}
\]
Its strict fixed surface is s=t zero, with the literal
source coordinates x,y. In the actual period block
the selected/ordinary comparison has index zero and
therefore no additional q power on these four sigma
charts. Its complete point factor is the same \(\psi\)
whose restriction was proved one. The source period has
line multiplier \(P_G^a=1\) on the central source.
Hence SS.8 closes across the entire selected source
disc in the specified reference, with no extra fibre
divisor or annular integer.

The original finite cap lies on the large-sigma domain.
It has no point on these sigma-zero strict surfaces.
Its retained \(AB^{j_\gamma}\) square is unchanged in
the surrounding comparison family. No cap character
has been used to alter SS.8.

\par

\medskip\noindent\textbf{SS.3. The exact strict surface resolution}

\par\nobreak

For each strict face, the four actual 238-chart tangent
cones project to the source rays
\((0,1),(1,1),(2,1),(3,1),(4,1)\) in
\((\operatorname{ord}\chi^4,\operatorname{ord}\chi\eta)\).
They give the same source morphism
\[
 x=a^{k+1}b^k,\quad y=a^{3-k}b^{4-k},\quad z=ab,
 \qquad xy=z^4,\quad0\le k\le3
 \tag{SS.10}
\]
at each O/T quotient point. Thus these strict-source
stars are the actual minimal A3 source resolution, by
their invariant maps and all four inverse charts.
The other two-torsion pair gives the actual A1 source
resolution \((\rho,z)\mapsto(\rho,\rho z,\rho z^2)\).
Away from those quotient points the fixed source maps
are SS.2, the identity source charts SS.9 and the native
I2 graph charts. They have inverse source coordinates
on each indicated open. They glue by the same source
functions. Therefore the completed strict source
surface in every index is the actual minimal native
surface \(S_0\), through this chartwise isomorphism.

The independent source comparison proves
\(p_1(\bar S)=-6\). The exact two A3 comparisons each
change p1 by \(-15/2\), and the actual A1 comparison
changes it by \(-3\). Thus the tangent number on each
completed strict surface is
\[
 \int_{S_0}p_1=-24.
 \tag{SS.11}
\]
The proof, with its invariant volume and common radial
boundary field, is
\path{part3_e6_selected_infinity_a3_source_working.md};
its independent peer checks are retained in
\path{fourfold_audit_e6_selected_infinity_a3_source.json}.

\par

\medskip\noindent\textbf{SS.4. Integral normal intersections and their boundary characters}

\par\nobreak

Write \(D_j=c_1(N_{+,j})-c_1(N_{-,j})\). Formula SS.8
and the literal section intersections SS.5 give
\[
 D_j=-(4-j)(O+T)+(2j+1)L,
\]
\[
 D_j^2=-\frac{(4-j)^2}{2}-(4-j)(2j+1).
 \tag{SS.12}
\]
The four orbifold squares are
\(-12,-27/2,-12,-15/2\).

The normal degrees in the three A3 exceptional curves,
in negative/positive physical order, are precisely
\[
\begin{array}{c|ccc}
 j&C_1&C_2&C_3\\\hline
 0&(1,0)&(0,0)&(0,0)\\
 1&(1,0)&(-1,0)&(1,0)\\
 2&(0,0)&(0,0)&(1,0)\\
 3&(0,0)&(0,0)&(0,0).
\end{array}
 \tag{SS.13}
\]
These are the normal monomial clutches of the actual
four-cones. The certificate recomputes each degree
using two independent auxiliary circles and retains
every matrix and unit exponent. Use the integral
intersection matrix
\[
 I=\begin{pmatrix}-2&1&0\\1&-2&1\\0&1&-2\end{pmatrix}.
\]
For each row let d be the positive-minus-negative
degree vector. The relative lift is the exact rational
solution \(Iq=d\), and the relative square is
\(d^tI^{-1}d\). This gives
\(-3/4,-1,-3/4,0\) at each O/T point. The boundary
class is retained: in the ordered source chain its
restriction is \(d_1+2d_2+3d_3\bmod4=-(j+1)\).
This is the actual line clutch carried by the
quotient coefficient character in SS.8. Thus the
comparison uses its specified flat boundary holonomy.

The sign of this boundary class is fixed by a literal
section map. The coefficient-character-one line has
the equivariant section \(\chi\). Its strict divisor
is the source ray \((4,1)\), meeting the third
exceptional curve once. Its boundary Chern class is
therefore \(e_3\) in \(\operatorname{coker}I\).
The three displayed integral relations give
\(e_3=3e_1=-e_1\bmod4\). Consequently fibre
coefficient character h maps to boundary Chern
class \(-h\) in this ordered chain convention.
SS.8 has D coefficient character \(j+1\), agreeing
exactly with the class \(-(j+1)\) computed from
SS.13. This also specifies the generator used in
the rational relative comparison.

More explicitly, the boundary is the actual lens space
given by the quotient of the source sphere. Its cellular
complex has differential \(d_2=[4]\) and other positive
differentials zero, so its first and second rational
cohomology vanish. The long exact sequence of the pair
therefore identifies the relative and absolute second
rational cohomology of this resolution neighbourhood.
The integral intersection map in their zero-section and
fibre-disc bases is I. The actual integral normal line
has degree vector d; its unique rational relative lift
has coefficient vector \(I^{-1}d\). Its pairing with
the same absolute class is precisely
\(d^tI^{-1}d\). The full normal monomials identify its
boundary line with the stated flat orbifold character,
so matching the two connections uses that exact relative
class. This explains both the rational coefficients and
the retained integral boundary torsion in the comparison.

At the A1 source patch the complete line morphism is
\(\ell=\chi a\), with its actual order-two character.
The integral matrix \([-2]\) and degree \(\pm1\)
give the relative D square \(-1/2\), as already
proved with the exact boundary map in the mu2 note.
Combining these specific relative extensions with
SS.12 yields the resolved squares
\[
 (D_0^2,D_1^2,D_2^2,D_3^2)=(-14,-16,-14,-8).
 \tag{SS.14}
\]

\par

\medskip\noindent\textbf{SS.5. The resulting strict-surface part of the boundary ledger}

\par\nobreak

For these actual normal weights \((+1,-1)\), the
positive-hyperplane residue is
\((p_1-3D^2)/24\). Equations SS.11 and SS.14 give
\[
 (R_u(S_0),R_u(S_1),R_u(S_2),R_u(S_3))
       =\left(\frac34,1,\frac34,0\right).
 \tag{SS.15}
\]
Their sum is \(5/2\). The four isolated points of the
same actual mu2 resolution add \(1/4\), giving
\(11/4\) for this strict-source-plus-mu2 part of the
positive-hyperplane ledger. Equivalently, retaining
the A1 orbifold source patches and using the proved
zero total mu2 correction gives the same \(11/4\).
This equality follows through the literal common
boundaries and their line characters.

The physical line is \(M=U^{-1}\). Its cubic and
Pontryagin-linear integrals both change sign, so the
same strict-source-plus-mu2 part has physical value
\(-11/4\). The O/T interior points and curves,
the exceptional-source surfaces, the refinement
surfaces and the moving selected reference cap
are separate entries in the existing boundary
equation. They have not been omitted or replaced
by the number in SS.15. No total return is assigned.

